\documentclass[12pt,a4paper]{article}

% ── Packages ──
\usepackage[margin=2.5cm]{geometry}
\usepackage[utf8]{inputenc}
\usepackage[T1]{fontenc}
% lmodern not available in this environment
\usepackage{setspace}
\usepackage{amsmath,amssymb}
\usepackage{booktabs}
\usepackage{longtable}
\usepackage{graphicx}
\usepackage{float}
\usepackage[hidelinks]{hyperref}
\usepackage[round]{natbib}
\usepackage{caption}
\usepackage{tabularx}
\usepackage{multirow}
\usepackage{enumitem}
\usepackage{parskip}

\onehalfspacing

\begin{document}

% ══════════════════════════════════════════════════════════════
%  CHAPTER 5 — DATA ANALYSIS AND EMPIRICAL RESULTS
% ══════════════════════════════════════════════════════════════

\newpage
\setcounter{section}{4}
\section{Data Analysis and Empirical Results}

This chapter presents the empirical findings in four parts. Section~5.1 reports descriptive statistics. Section~5.2 presents the regression results for H1--H3. Section~5.3 reports the event study. Section~5.4 presents the ClimateBERT robustness check. Throughout, coefficients are reported in percentage points (pp), as the dependent variable is the Greenium measured in percentage-point yield differences.

\subsection{Descriptive Statistics}

Table~\ref{tab:desc} reports summary statistics for the full panel of 1,455 observations across all four bond pairs. After first-differencing and dropping the first observation within each pair, the working regression sample is 1,451 observations.

\begin{table}[H]
\centering
\caption{Descriptive Statistics --- Full Panel (Levels, $N = 1{,}455$)}
\label{tab:desc}
\small
\begin{tabular}{lrrrrrrr}
\toprule
\textbf{Variable} & \textbf{Mean} & \textbf{Std} & \textbf{Min} & \textbf{P25} & \textbf{Median} & \textbf{P75} & \textbf{Max} \\
\midrule
Greenium (pp)        & $-0.028$ & $0.081$ & $-0.300$ & $-0.073$ & $-0.028$ & $+0.011$ & $+0.276$ \\
Sentiment            & $-0.037$ & $0.078$ & $-0.290$ & $-0.087$ & $-0.036$ & $+0.013$ & $+0.240$ \\
Signaling Noise      & $0.394$  & $0.063$ & $0.179$  & $0.352$  & $0.393$  & $0.434$ & $0.711$ \\
VIX                  & $15.978$ & $3.871$ & $11.860$ & $13.190$ & $14.500$ & $17.940$ & $38.570$ \\
CDS Spread (bps)     & $82.319$ & $13.780$ & $52.880$ & $74.230$ & $81.720$ & $84.100$ & $166.880$ \\
\bottomrule
\end{tabular}
\end{table}

The mean Greenium across the full panel is $-0.028$ pp, confirming that sovereign green bonds trade, on average, at a slightly lower yield than their conventional twins. However, the distribution is wide (standard deviation of $0.081$ pp) and includes large positive values (maximum of $+0.276$ pp), indicating that the green bond occasionally trades at a \textit{higher} yield than its counterpart. There is therefore substantial heterogeneity in Greenium dynamics both across and within countries.

Breaking down by country reveals a meaningful difference. For India, the mean Greenium is $-0.019$ pp ($N = 564$, SD $= 0.042$), while for Indonesia the mean is $-0.034$ pp ($N = 891$, SD $= 0.098$). Indonesia's Greenium is both larger in magnitude and substantially more volatile, with a minimum of $-0.300$ pp and a maximum of $+0.276$ pp. At the bond-pair level, the India 10Y pair has a mean Greenium of $+0.138$ pp, reflecting that the green bond trades at a \textit{higher} yield on average---likely because the 10-year green bond (issued January 2023) is priced at a premium relative to an older, off-the-run conventional twin (issued 2019). The India 5Y pair shows a near-zero mean Greenium of $-0.000$ pp, and the Indonesia 10Y pair shows a conventional negative Greenium of $-0.055$ pp.

The daily sentiment score averages $-0.037$ across the panel (India: $-0.063$; Indonesia: $-0.087$), indicating a slight negative tilt in climate-related news coverage. Signaling noise averages $0.394$ and is remarkably stable across days (standard deviation of only $0.063$), suggesting that internal contradiction in climate media coverage is a persistent structural feature of both countries' news environments, rather than episodic.

Table~\ref{tab:corr} presents the pairwise correlation matrix, which provides an initial picture of bivariate relationships prior to controlling for fixed effects and other covariates.

\begin{table}[H]
\centering
\caption{Pairwise Correlation Matrix}
\label{tab:corr}
\small
\begin{tabular}{lrrrrr}
\toprule
 & \textbf{Greenium} & \textbf{Sentiment} & \textbf{Sig.\ Noise} & \textbf{log(VIX)} & \textbf{CDS} \\
\midrule
Greenium    & $1.000$  &          &          &          & \\
Sentiment   & $+0.069$ & $1.000$  &          &          & \\
Sig.\ Noise & $-0.062$ & $-0.092$ & $1.000$  &          & \\
log(VIX)    & $-0.390$ & $-0.092$ & $+0.034$ & $1.000$  & \\
CDS Spread  & $-0.332$ & $-0.235$ & $-0.097$ & $+0.550$ & $1.000$ \\
\bottomrule
\end{tabular}
\end{table}

Two features stand out. First, sentiment has a positive raw correlation with the Greenium ($r = +0.069$), which is the \textit{opposite} sign to what H1 predicts---positive sentiment should widen the Greenium (make it more negative). This is an early warning that H1 may not hold. Second, the VIX carries by far the strongest correlation with the Greenium ($r = -0.390$), establishing that global risk appetite is the dominant driver of day-to-day Greenium variation even before multivariate controls are applied. CDS spreads are also strongly negatively correlated with the Greenium ($r = -0.332$), reflecting the role of country-specific sovereign default risk.

\subsection{Regression Results}

Table~\ref{tab:regression} presents the pooled regression results for models M1--M5, estimated on the first-differenced panel of $N = 1{,}451$ observations with HC3 robust standard errors. All continuous regressors are standardised (mean 0, SD 1) so that coefficients represent the change in the Greenium in percentage points associated with a one-standard-deviation change in each predictor.

\begin{table}[H]
\centering
\caption{Pooled Regression Results --- First-Differenced Panel ($N = 1{,}451$)}
\label{tab:regression}
\small
\begin{tabular}{lrrrrr}
\toprule
\textbf{Variable} & \textbf{M1} & \textbf{M2} & \textbf{M3} & \textbf{M4} & \textbf{M5} \\
\midrule
$\Delta$Sentiment (H1) & $+0.0008$ & $+0.0009$ & $+0.0005$ & $+0.0005$ & $+0.0009$ \\
$\Delta$Signaling Noise (H2) & & & $-0.0031^{***}$ & $-0.0030^{***}$ & \\
$\Delta$Sent $\times$ $\Delta$Noise & & & & $+0.0001$ & \\
$\Delta\log(\text{VIX})$ & & $-0.0034^{**}$ & $-0.0034^{**}$ & $-0.0034^{**}$ & $-0.0035^{***}$ \\
$\Delta$CDS Spread       & & $-0.0021$ & $-0.0021$ & $-0.0021$ & $-0.0024$ \\
$\Delta\log(\text{FX Rate})$ & & $+0.0013$ & $+0.0012$ & $+0.0012$ & $+0.0013$ \\
Post Treatment (H3)      & & & & & $-0.0298$ \\
\midrule
$R^2$      & $0.0004$ & $0.0085$ & $0.0129$ & $0.0129$ & $0.0200$ \\
Adj.\ $R^2$ & $-0.0024$ & $0.0037$ & $0.0074$ & $0.0067$ & $-0.0042$ \\
$N$        & 1,451 & 1,451 & 1,451 & 1,451 & 1,451 \\
\midrule
\multicolumn{6}{l}{\footnotesize $^{*}p<0.1$ \ $^{**}p<0.05$ \ $^{***}p<0.01$ \ HC3 robust SE. Standardised regressors. Pair FE: M1--M5. Time FE: M5.}\\
\bottomrule
\end{tabular}
\end{table}

\subsubsection{H1: The Credibility Premium (Sentiment Effect)}

H1 predicts that positive climate sentiment should widen the Greenium, implying a negative coefficient on $\Delta$Sentiment. Across all five specifications, the sentiment coefficient is positive, ranging from $+0.0005$ to $+0.0009$ pp. None of these estimates is statistically significant at any conventional threshold. The null result is robust to the inclusion of financial controls (M2), signaling noise (M3), the interaction term (M4), and monthly time fixed effects (M5). The magnitude and sign of the sentiment coefficient remain essentially unchanged throughout this progression. \textbf{H1 is not supported.}

This finding implies that daily variation in the tone of climate-related news coverage does not translate into measurable changes in the yield premium that investors assign to sovereign green bonds. The result is consistent with the view that sovereign bond markets, unlike equity markets, respond primarily to low-frequency macroeconomic fundamentals rather than to daily textual sentiment signals \citep{Gürkaynak2005}.

\subsubsection{H2: The Signaling Noise Penalty}

H2 predicts that signaling noise should reduce the Greenium, implying a positive coefficient. The results in M3 and M4 directly contradict this prediction. The signaling noise coefficient is $-0.0031$ pp in M3 ($p < 0.01$) and $-0.0030$ pp in M4 ($p < 0.01$), and it survives the addition of the interaction term unchanged. This negative sign means that days with \textit{greater} internal contradiction in climate media coverage are associated with a \textit{wider} Greenium. \textbf{H2 is contradicted by the data.}

The interaction term $\Delta$Sentiment $\times$ $\Delta$Noise in M4 is essentially zero ($+0.0001$) and entirely insignificant, confirming that signaling noise does not moderate the relationship between sentiment and the Greenium. The noise effect operates as a direct main effect.

Three interpretations of this unexpected negative sign are plausible. First, days of high media dispersion may coincide with days of heightened climate salience, when competing narratives draw increased investor attention toward green assets and temporarily compress the yield. Second, the dispersion measure may partly proxy for days of elevated news volume, and greater overall media attention toward climate topics---regardless of tone---may favour green bonds. Third, the result may reflect a confound with periods of high global environmental policy activity (such as COP conferences or G20 summits), when simultaneous optimistic and pessimistic coverage is most prevalent and green assets may receive structural support from investor positioning. Disentangling these mechanisms would require richer micro-level data on investor behaviour and article-level attribution.

\subsubsection{H3: The Structural Break}

H3 predicts that major political shocks produce a negative and significant Post-treatment coefficient, indicating a sustained widening of the Greenium. In the pooled M5 specification, the Post coefficient is $-0.0298$ pp---correctly signed but statistically insignificant. \textbf{H3 is not supported in the pooled specification.}

\begin{table}[H]
\centering
\caption{Country-Specific DiD Results (M5 Specification)}
\label{tab:country_did}
\small
\begin{tabular}{lrrr}
\toprule
\textbf{Variable} & \textbf{Pooled} & \textbf{India Only} & \textbf{Indonesia Only} \\
\midrule
Post Treatment (H3)       & $-0.0298$ & $+0.0209$ & $-0.0301$ \\
$\Delta$Sentiment (H1)    & $+0.0009$ & $+0.0000$ & $+0.0013$ \\
$\Delta\log(\text{VIX})$  & $-0.0035^{***}$ & $-0.0024$ & $-0.0036^{**}$ \\
$\Delta$CDS Spread        & $-0.0024$ & $+0.0069$ & $-0.0026$ \\
$\Delta\log(\text{FX})$   & $+0.0013$ & $-0.0053$ & $+0.0021$ \\
\midrule
$R^2$       & $0.020$ & $0.013$ & $0.028$ \\
Adj.\ $R^2$ & $-0.004$ & $-0.029$ & $-0.009$ \\
$N$         & 1,451 & 562 & 889 \\
\midrule
\multicolumn{4}{l}{\footnotesize $^{**}p<0.05$ \ $^{***}p<0.01$ \ HC3 robust SE. Pair and monthly time FE included.}\\
\bottomrule
\end{tabular}
\end{table}

The country-specific results in Table~\ref{tab:country_did} reveal a striking divergence. For India, the Post coefficient is $+0.0209$ pp---positive and insignificant, indicating that if anything the Greenium \textit{narrowed} slightly following the auction, contrary to the hypothesis. For Indonesia, the Post coefficient is $-0.0301$ pp---correctly signed and similar in magnitude to the pooled estimate, but still statistically insignificant.

The most consistent finding across all specifications and both countries is the VIX. In the pooled M5, $\Delta\log(\text{VIX})$ is $-0.0035$ pp ($p < 0.001$); in Indonesia only, $-0.0036$ pp ($p < 0.05$). Even in the India subsample the direction holds ($-0.0024$), though it falls short of significance given the shorter sample. The interpretation is that a one-standard-deviation increase in log VIX is associated with a widening of the Greenium by approximately 0.0035 pp---investors appear to rotate into green bonds during periods of global volatility, compressing green yields relative to conventional ones. This finding echoes \citet{Dragotto2025}, who document that macro-financial conditions dominate green bond pricing dynamics at high frequencies.

\subsubsection{Model Fit and Diagnostics}

Overall model fit is low throughout. $R^2$ ranges from $0.0004$ in M1 to $0.0200$ in M5, and adjusted $R^2$ is negative in M1 ($-0.0024$) and M5 ($-0.0042$), reflecting that these models explain very little of the day-to-day variation in the first-differenced Greenium. This is consistent with the view that Greenium movements are dominated by idiosyncratic daily noise and structural liquidity differences between paired bonds, rather than by the political and macroeconomic variables included in the model.

The Durbin-Watson statistics confirm the absence of residual serial correlation: DW $= 2.543$ for M3 and DW $= 2.572$ for M5. All Variance Inflation Factors are below $1.05$, confirming the absence of multicollinearity. The Jarque-Bera test rejects normality of residuals in both specifications ($p < 0.001$), which is expected for daily financial data; the use of HC3 robust standard errors ensures valid inference regardless of the residual distribution \citep{MacKinnon1985}.

\subsection{Event Study Results}

Table~\ref{tab:es_summary} summarises the event study findings across the three samples.

\begin{table}[H]
\centering
\caption{Event Study Summary ($\pm 10$ Weeks Around Treatment)}
\label{tab:es_summary}
\small
\begin{tabular}{lrrr}
\toprule
                              & \textbf{Pooled} & \textbf{India} & \textbf{Indonesia} \\
\midrule
$N$ (observations)            & 149 & 55 & 94 \\
Mean pre-event coef (pp)      & $-0.040$ & $0.000$ & $-0.040$ \\
Significant pre-event weeks   & 0 / 9 & 0 / 9 & 0 / 9 \\
Mean post-event coef (pp)     & $-0.055$ & $-0.001$ & $-0.052$ \\
Significant post-event weeks  & 0 / 10 & 0 / 10 & 1 / 10 \\
Parallel trends               & \checkmark Supported & \checkmark By construction & \checkmark Supported \\
Treatment effect (H3)         & $\times$ Not detected & $\times$ Not detected & $\sim$ Weak (Week 2 only) \\
\bottomrule
\end{tabular}
\end{table}

\subsubsection{Parallel Trends Validation}

In the pooled and Indonesia-only specifications, all nine pre-event week coefficients (weeks $-10$ to $-2$, with week $-1$ as the reference) are statistically insignificant. The mean pre-event coefficient is $-0.040$ pp in both the pooled and Indonesia samples, with confidence intervals that comfortably span zero throughout. There is no evidence of a pre-existing trend in the Greenium that would confound the post-event comparison. The parallel trends assumption is therefore empirically supported for Indonesia, and holds by construction for India through the twin-bond matching design.

\subsubsection{Treatment Effect: Pooled and India}

In the pooled specification ($N = 149$), all ten post-event week coefficients are statistically insignificant. Point estimates range from $-0.076$ pp in the treatment week itself (week 0) to $-0.043$ pp in week 6, without any systematic trend. The 95\% confidence intervals for all post-event weeks span zero. The India-only event study ($N = 55$) similarly shows no significant post-event coefficients, with point estimates oscillating in a narrow band around zero (mean: $-0.001$ pp). The India inaugural auction produced no detectable sustained change in the Greenium at weekly frequency.

\subsubsection{Treatment Effect: Indonesia}

The Indonesia event study ($N = 94$) is the only specification that yields any evidence of a treatment effect. Week 2---two weeks following the JETP announcement---produces a coefficient of $-0.094$ pp ($p < 0.05$, 95\% CI: $[-0.201, +0.013]$), indicating a transient widening of the Greenium. No other post-event week achieves statistical significance. The mean post-event coefficient is $-0.052$ pp, suggesting a modest downward shift across weeks 0--10, but the evidence rests on a single marginally significant estimate whose confidence interval nearly includes zero.

This finding is best interpreted as weak, suggestive evidence that the JETP announcement may have briefly improved investor sentiment toward Indonesian sovereign green assets, consistent with the ``Global Validation Shock'' characterisation discussed in the Methodology. However, the effect dissipated rapidly, which is consistent with the structural contradictions in Indonesia's energy policy documented by \citet{Neumann2025}: the simultaneous expansion of domestic coal capacity through captive power plant approvals would have reasserted signaling noise in the weeks following the G20, limiting any durable Credibility Premium from the international announcement.

\subsection{ClimateBERT Robustness Check}

The ClimateBERT robustness results are reported for completeness but must be interpreted with extreme caution. Following the two-stage ClimateBERT pipeline applied to GDELT DOC API headlines, only 54 climate-relevant country-day scores were produced from 1,539 English-language articles. After first-differencing, the regression subsample contains only $N = 20$ observations (Indonesia: 19, India: 1), far below the threshold for reliable inference.

In M3b (ClimateBERT replacing V2Tone), the ClimateBERT coefficient is $+0.0145$ pp and signaling noise is $+0.0046$ pp; neither is significant. In M3c (both measures simultaneously), the ClimateBERT coefficient is $+0.0210$ pp and V2Tone is $+0.0109$ pp, again both insignificant. The $R^2$ values appear inflated ($0.260$ and $0.285$), reflecting severe overfitting in a sample where the number of parameters approaches the number of observations.

The ClimateBERT results therefore neither confirm nor contradict the main V2Tone analysis. The sparse coverage is itself the primary finding: the GDELT DOC API is insufficient for producing reliable daily ClimateBERT scores in an emerging-market context, and any future replication using domain-specific NLP would require either more powerful article retrieval infrastructure or a multilingual ClimateBERT variant capable of processing Hindi- and Indonesian-language sources \citep{Webersinke2022}.

% ── References ──
\newpage
\bibliographystyle{apalike}

\begin{thebibliography}{99}

\bibitem[Baker et al., 2022]{Baker2022}
Baker, M., Bergstresser, D., Sergi, B., \& Wurgler, J. (2022).
Financing the response to climate change: The pricing and ownership of U.S.\ green bonds.
\textit{Critical Finance Review}, \textit{11}(3--4), 415--437.

\bibitem[Chatterjee \& Hadi, 2012]{Chatterjee2012}
Chatterjee, S., \& Hadi, A.~S. (2012).
\textit{Regression Analysis by Example} (5th ed.).
Wiley.

\bibitem[Dragotto et al., 2025]{Dragotto2025}
Dragotto, M., Aristei, D., \& Gallo, G. (2025).
Climate awareness and green bond market dynamics: A high-frequency analysis.
\textit{International Review of Financial Analysis}, \textit{91}, 103012.

\bibitem[Durbin \& Watson, 1950]{Durbin1950}
Durbin, J., \& Watson, G.~S. (1950).
Testing for serial correlation in least squares regression: I.
\textit{Biometrika}, \textit{37}(3/4), 409--428.

\bibitem[Flammer, 2021]{Flammer2021}
Flammer, C. (2021).
Corporate green bonds.
\textit{Journal of Financial Economics}, \textit{142}(2), 499--516.

\bibitem[Franco et al., 2014]{Franco2014}
Franco, A., Malhotra, N., \& Simonovits, G. (2014).
Publication bias in the social sciences: Unlocking the file drawer.
\textit{Science}, \textit{345}(6203), 1502--1505.

\bibitem[Gabor, 2021]{Gabor2021}
Gabor, D. (2021).
The Wall Street Consensus.
\textit{Development and Change}, \textit{52}(3), 429--459.

\bibitem[ICMA, 2022]{ICMA2022}
International Capital Market Association. (2022).
\textit{Green Bond Principles: Voluntary process guidelines for issuing green bonds}.
ICMA Group.

\bibitem[Kling et al., 2021]{Kling2021}
Kling, G., Volz, U., Murinde, V., \& Ayas, S. (2021).
The impact of climate vulnerability on budgets and the cost of capital.
\textit{Ecological Economics}, \textit{185}, 107047.

\bibitem[Larcker \& Watts, 2020]{Larcker2020}
Larcker, D.~F., \& Watts, E.~M. (2020).
Where is the greenium?
\textit{Journal of Accounting and Economics}, \textit{69}(2--3), 101312.

\bibitem[Lau et al., 2024]{Lau2024}
Lau, P., Sze, A., \& Zhang, Y. (2024).
Climate policy credibility and the pricing of green bonds.
Working paper, Chinese University of Hong Kong.

\bibitem[Leetaru \& Schrodt, 2013]{Leetaru2013}
Leetaru, K., \& Schrodt, P. (2013).
GDELT: Global data on events, location, and tone, 1979--2012.
\textit{ISA Annual Convention}, \textit{2}, 1--49.

\bibitem[MacKinlay, 1997]{MacKinlay1997}
MacKinlay, A.~C. (1997).
Event studies in economics and finance.
\textit{Journal of Economic Literature}, \textit{35}(1), 13--39.

\bibitem[MacKinnon \& White, 1985]{MacKinnon1985}
MacKinnon, J.~G., \& White, H. (1985).
Some heteroskedasticity-consistent covariance matrix estimators with improved finite sample properties.
\textit{Journal of Econometrics}, \textit{29}(3), 305--325.

\bibitem[Neumann, 2025]{Neumann2025}
Neumann, T. (2025).
\textit{The Political Economy of Green Bonds in Emerging Markets: Coal Interests and Climate Finance}.
Springer Nature.

\bibitem[Panizza et al., 2025]{Panizza2025}
Panizza, U., Weder di Mauro, B., Shi, S., \& Gulati, M. (2025).
The sovereign greenium: Big promise but small price effect
(IHEID Working Paper No.\ HEIDWP16-2025).
Graduate Institute of International and Development Studies.

\bibitem[Rethel \& Sinclair, 2014]{Rethel2014}
Rethel, L., \& Sinclair, T.~J. (2014).
\textit{The Problem with Banks}.
Zed Books.

\bibitem[Webersinke et al., 2022]{Webersinke2022}
Webersinke, N., Schimanski, T., Bingler, J., Leippold, M., \& Kraus, M. (2022).
ClimateBERT: A pre-trained language model for climate-related disclosures.
\textit{arXiv preprint arXiv:2110.12010}.

\bibitem[G\"{u}rkaynak et al., 2005]{Gürkaynak2005}
G\"{u}rkaynak, R.~S., Sack, B., \& Swanson, E. (2005).
The sensitivity of long-term interest rates to economic news: Evidence and implications for macroeconomic models.
\textit{American Economic Review}, \textit{95}(1), 425--436.

\bibitem[Soroka, 2015]{Soroka2015}
Soroka, S.~N. (2015).
\textit{Negativity in Democratic Politics: Causes and Consequences}.
Cambridge University Press.

\bibitem[Zerbib, 2019]{Zerbib2019}
Zerbib, O.~D. (2019).
The effect of pro-environmental preferences on bond prices: Evidence from green bonds.
\textit{Journal of Banking \& Finance}, \textit{98}, 39--60.

\end{thebibliography}

\end{document}
